\documentclass{article} % Definisi jenis dokumen

%%%%% Definisi paket-paket yang seharusnya digunakan %%%%%
\usepackage[utf8]{inputenc} % paket encoding input utf8
\usepackage[T1]{fontenc} % paket encoding huruf latin

%%%%% Definisi paket-paket yang digunakan sesuai kebutuhan %%%%%
\usepackage[yyyymmdd,hhmmss]{datetime} % paket tanggal-waktu
\usepackage{graphicx} % paket grafik/gambar
\usepackage[english]{babel} % paket modifikasi label/caption pada bahasa tertentu
\usepackage{geometry} % paket ukuran kertas dan margin
\usepackage{tikz} % paket menggambar di latex

%%%%% Pengaturan ukuran kertas dan margin %%%%%
\geometry{
	a4paper,
	left=10mm,
	right=10mm,
	top=15mm,
	bottom=15mm,
}

%%%%% Pengaturan perintah informasi perangkat lunak (hanya untuk GNU/Linux) %%%%%
\newcommand{\ShowOsVersion}{
	\immediate\write18{\unexpanded{foo=`uname -sro` && echo "${foo}" > tmp.tex}}
	\input{tmp}\immediate\write18{rm tmp.tex}
}

\newcommand{\ShowTexVersion}{
	\immediate\write18{\unexpanded{foo=`pdflatex -version | head -n1 | cut -d' ' -f1,2` && echo "${foo}" > tmp.tex}}
	\input{tmp}\immediate\write18{rm tmp.tex}
}

\usetikzlibrary{positioning}
\newdimen\nodeDist
\nodeDist=35mm

\begin{document}
	%%%%%%%%%%%%%%%%%%%%%%%%%%%%%%%%%%%%%%%%%%%%%%%%%%%%%%%%%%%%%%%%%
	
	\begin{titlepage}
		
		\centering % untuk membuat tengah teks
		
		{
			\LARGE % pakai font besar
			\bf % pakai font BOLD
			VIBIO Design Plan 
		}
		
		\bigskip
		{\Large \bf Achmadi ST MT}
		\vfill % menambahkan ruang kosong vertikal
		
		\raggedright
		\noindent Buku ini ditulis dengan:\\ % tanda \\ menambahkan garis baru
		OS : \ShowOsVersion \\
		TeX : \ShowTexVersion \\
		Update: {\today} at \currenttime\\
	\end{titlepage}
	
	%%%%%%%%%%%%%%%%%%%%%%%%%%%%%%%%%%%%%%%%%%%%%%%%%%%%%%%%%%%%%%%%%
	
	\section{Considerations}
	
	Here below a summary of things to considered while chosing this design plan:
	
	\begin{itemize}
		\item Physical interface requirement, e.g: Buttons, LEDs, Knobs, etc
		
		\item Speech Recognition implementation, the choice is either None, Offline, or Online
		
		\item Custom built prototype requirement, either can be custom built or just by prebuilt tablet or similar
		
		\item Terapist's interaction requirement, either as guided thereapist or none
	\end{itemize}

	\section{Planning Diagram}
	
	\begin{tikzpicture}[
		node/.style={%
			draw,
			rectangle,
		},
		]
		
		\node [node] (A) {VIBIO Hardware};
		\path (A) ++(-135:\nodeDist) node [node] (B) {Tablet};
		\path (A) ++(-45:\nodeDist) node [node] (C) {RaspberryPi};
		\path (C) ++(-135:\nodeDist) node [node] (D) {LCD Only};
		\path (C) ++(-45:\nodeDist) node [node] (E) {ButtonLeds PCB};
		
		\draw (A) -- (B) node [left,pos=0.25] {Prebuilt}(A);
		\draw (A) -- (C) node [right,pos=0.25] {Custom}(A);
		\draw (C) -- (D) node [left,pos=0.25] {Touch Only}(A);
		\draw (C) -- (E) node [right,pos=0.25] {Physical Iface}(A);
	\end{tikzpicture}

	\vfill % menambahkan ruang kosong vertikal

	\begin{tikzpicture}[
		node/.style={%
			draw,
			rectangle,
		},
		]
		
		\node [node] (A) {VIBIO SpeecRecognition};
		\path (A) ++(-135:\nodeDist) node [node] (B) {Terapist};
		\path (A) ++(-90:\nodeDist) node [node] (C) {Python/C/C++};
		\path (A) ++(-45:\nodeDist) node [node] (D) {Webserver};
		\path (B) ++(-135:\nodeDist) node [node] (E) {Parents};
		\path (B) ++(-90:\nodeDist) node [node] (F) {Expert};
		\path (C) ++(-90:\nodeDist) node [node] (G) {HTK/HMM ?};
		\path (D) ++(-45:\nodeDist) node [node] (H) {Tanya Pak Bagus};
		
		\draw (A) -- (B) node [left,pos=0.25] {None}(A);
		\draw (A) -- (C) node [below,pos=0.25] {Offline}(A);
		\draw (A) -- (D) node [right,pos=0.25] {Online}(A);
		\draw (B) -- (E) node [left,pos=0.25] {Guided}(B);
		\draw (B) -- (F) node [below,pos=0.25] {Trained}(B);
		\draw (C) -- (G) node [below,pos=0.25] {Precompiled}(C);
		\draw (D) -- (H) node [right,pos=0.25] {Decent-Internet}(D);
	\end{tikzpicture}

	\vfill % menambahkan ruang kosong vertikal

	\begin{tikzpicture}[
		node/.style={%
			draw,
			rectangle,
		},
		]
		
		\node [node] (A) {VIBIO Operation};
		\path (A) ++(-135:\nodeDist) node [node] (B) {Solo};
		\path (A) ++(-45:\nodeDist) node [node] (C) {Terapist};
		\path (C) ++(-135:\nodeDist) node [node] (D) {Parents};
		\path (C) ++(-45:\nodeDist) node [node] (E) {Expert};
		\path (B) ++(-135:\nodeDist) node [node] (F) {How?};
		
		\draw (A) -- (B) node [left,pos=0.25] {Full Auto}(A);
		\draw (A) -- (C) node [right,pos=0.25] {Interactions}(A);
		\draw (C) -- (D) node [left,pos=0.25] {Guided}(A);
		\draw (C) -- (E) node [right,pos=0.25] {Trained}(A);
		\draw (B) -- (F) node [right,pos=0.25] {Used by Toddler}(B);
	\end{tikzpicture}
	
	\newpage
	\section{Current Stage of Development}
	
	\subsection{Main Program}
	
	This main program written as a web interface packed in Electron (Nextron) and ready to be deployed on any capable computer (either x86 or armv7).
	
	Below is the example of the program running on top of Xorg in a Arch Linux operating system.
	
	\begin{figure}[!ht]                                  % [!ht] sebagai floatifier really here-top
		\centering                                       % Menengahkan gambar
		
		% Tampilkan dalam ukuran lebar 200pt gambar image/organik_bro
		% Pastikan berkas gambar atau foldernya di folder yang sama dengan berkas tex nya.
		\includegraphics[width=350pt]{images/vibio}
		
		\caption{Main Program}                         % Label gambar
	\end{figure}

	\subsection{ARM Porting}
	
	As Arch Linux also available for ARM devices such as RaspberryPi, the main program already tested in a RaspberryPi 3B unit
	to fulfill the demand of custom built hardware if required
	
	\begin{figure}[!ht]                                  % [!ht] sebagai floatifier really here-top
		\centering                                       % Menengahkan gambar
		
		% Tampilkan dalam ukuran lebar 200pt gambar image/organik_bro
		% Pastikan berkas gambar atau foldernya di folder yang sama dengan berkas tex nya.
		\includegraphics[width=350pt]{images/on_raspi}
		
		\caption{Run on Raspi-3B}                         % Label gambar
	\end{figure}
\end{document}